\section{Genuine finite Laurent coefficients of the entire scalar-extended module}

\begin{theorem}\label{thm:laurent-module-actual-coefficients}
\lean{EnvelopeProperty.LaurentModuleExpansion.coefficients, EnvelopeProperty.LaurentModuleExpansion.coefficients_tmul, EnvelopeProperty.LaurentModuleExpansion.coefficients_T_tmul, EnvelopeProperty.LaurentModuleExpansion.coefficients_symm, EnvelopeProperty.LaurentModuleExpansion.eq_sum_T_tmul}
\leanok
For every original commutative ring $R$ and original $R$-module $M$,
there is a canonical $R$-linear equivalence
\[
 R[T,T^{-1}]\otimes_R M\simeq\bigoplus_{i\in\mathbf Z}M.
\]
The coefficient of $F\otimes m$ at $i$ is $F_i m$. Every original
tensor is exactly the finite sum of $T^i\otimes m_i$ for its actual
coefficients, and these coefficients determine it uniquely. No
finiteness or freeness of $M$ is required.
\uses{def:native-multiplicative-group}
\end{theorem}
\begin{proof}
The original Laurent algebra has its genuine monomial basis indexed
by all integers. Tensoring this actual basis gives the canonical
finitely supported coefficient equivalence. Its pure-tensor formula
is the actual basis-coordinate formula. The inverse equivalence is
the finite sum of the original basis monomials tensored with their
module coefficients, so its inverse identities give reconstruction
and uniqueness, including the zero tensor.
\end{proof}

\begin{theorem}\label{thm:laurent-module-actual-evaluation}
\lean{EnvelopeProperty.LaurentModuleExpansion.evaluate, EnvelopeProperty.LaurentModuleExpansion.evaluate_coeff, EnvelopeProperty.LaurentModuleExpansion.evaluate_T_tmul}
\leanok
Let $S$ be an original commutative ring with an $R$-module structure,
and let $\phi:M\to S$ be an actual $R$-linear map. Applying $\phi$
to the actual module coefficients constructs a canonical whole
linear map into the original Laurent algebra $S[T,T^{-1}]$.
Its $i$th coefficient is exactly $\phi(m_i)$, and it sends each
$T^i\otimes m$ to $\phi(m)T^i$.
\uses{thm:laurent-module-actual-coefficients}
\end{theorem}
\begin{proof}
Apply the actual linear map coefficientwise on the finitely supported
module and then the inverse genuine Laurent coefficient equivalence.
The actual coefficient and monomial formulas follow from these two
inverse identities and the original pure-tensor coefficient formula.
Coefficients killed by $\phi$ are retained in the original expansion.
\end{proof}

\begin{theorem}\label{thm:laurent-module-actual-gauss-maximum}
\lean{EnvelopeProperty.LaurentModuleExpansion.laurentGauss_evaluate}
\leanok
For any valuation $w$ of $S$ and any $r>0$, the actual Laurent Gauss
value of the evaluated tensor is
\[
 \max_{i\in\operatorname{supp}(m_\bullet)}w(\phi(m_i))r^i.
\]
The support is that of the original module coefficients, so all
original negative exponents and coefficients killed by $\phi$ are
included. The empty maximum for the zero tensor is zero.
\uses{thm:laurent-module-actual-evaluation,spectrum:actual-laurent-gauss-valuation}
\end{theorem}
\begin{proof}
The genuine Laurent Gauss coefficient formula first gives the
maximum over the support of the evaluated Laurent polynomial.
The exact coefficient identity shows that this support lies in the
original module support. Enlarging the finite maximum by the
additional zero terms does not change it. Conversely every original
coefficient term is bounded by the actual Gauss value, including
terms whose image is zero. These two inequalities prove the formula.
\end{proof}

\begin{theorem}\label{thm:laurent-module-actual-finite-family}
\lean{EnvelopeProperty.LaurentModuleExpansion.terms, EnvelopeProperty.LaurentModuleExpansion.laurentGauss_finite_family}
\leanok
Every finite original family of scalar-extended module elements
produces one finite family of original pairs $(m_i,i)$, chosen before
$\phi$, $w$ and $r$. At every actual such evaluation, the maximum of
the original Laurent Gauss values is exactly the weighted maximum
$\max w(\phi(m_i))r^i$ over this same finite pair family.
\uses{thm:laurent-module-actual-gauss-maximum}
\end{theorem}
\begin{proof}
Take the finite union of the actual coefficient pairs of each
original tensor. The previous whole-tensor formula applies to each
family member. A finite maximum over this finite union is exactly
the iterated maximum over the original family and its original
coefficient supports. Repeated pairs and zero evaluation values do
not affect the maximum. This is an algebraic original-module
producer; identification of a genuine geometric global-section
space with this scalar extension remains the actual flat H0 step.
\end{proof}
